\chapter{Thiết kế hệ thống}

\begin{quote}
\itshape

Chương này tập trung trả lời câu hỏi "Hệ thống sẽ được xây dựng như thế nào?", chuyển đổi các yêu cầu từ Chương 2 thành một bản thiết kế kỹ thuật cụ thể, bao gồm kiến trúc tổng thể, cấu trúc dữ liệu, luồng xử lý logic và thiết kế giao diện người dùng.
\end{quote}

\section{Kiến trúc hệ thống tổng quan}

{\itshape
Trình bày Sơ đồ Kiến trúc (Architecture Diagram) cấp cao (ví dụ: Client-Server, Microservices) thể hiện các khối thành phần chính (ví dụ: Mobile App/Frontend, Backend/API Server, Database, Dịch vụ bên thứ ba như Google Maps API) và các luồng dữ liệu chính kết nối chúng.
}

\section{Thiết kế chi tiết}

\subsection{Sơ đồ lớp (Class Diagram)}

{\itshape
Trình bày Sơ đồ Lớp (Class Diagram) đại diện cho cấu trúc dữ liệu tĩnh của hệ thống, tập trung vào các thực thể cốt lõi (ví dụ: \texttt{User}, \texttt{Trip}, \texttt{Destination}, \texttt{Booking}, \texttt{Review}), bao gồm các thuộc tính quan trọng và mối quan hệ (ví dụ: 1-n, n-n) giữa chúng.
}

\subsection{Sơ đồ tuần tự (Sequence Diagram)}

{\itshape
Trình bày Sơ đồ Tuần tự (Sequence Diagram) để mô tả chi tiết luồng tương tác giữa các đối tượng/thành phần khi thực hiện một use case nghiệp vụ quan trọng, ví dụ như "Người dùng tạo kế hoạch chuyến đi" hoặc "Hệ thống gợi ý địa điểm".
}

\pagebreak

\section{Thiết kế giao diện}

\subsection{Các màn hình chính}

{\itshape
Trình bày các bản thiết kế (Figma mockups) của các giao diện người dùng (UI) then chốt của hệ thống (ví dụ: Màn hình tìm kiếm, Màn hình chi tiết địa điểm, Màn hình quản lý lịch trình).
}
